\section{Strengths and Weaknesses}

\subsection{Strengths}

% 每条挂具体机理，不写"模型简单""精度较高"这类空话。
\begin{enumerate}
  \item The force decomposition is explicit: disturbance, spring--damper constraint and
        actuator resultant force are derived separately, so each term can be checked
        against a physical measurement.
  \item The three problems share one equation of motion; only the decision variable
        changes. A reader who understands Problem 1 can follow Problem 3 without new
        mechanics.
  \item Every reported index is accompanied by a feasibility check on the angular
        velocity and acceleration limits, so the schedule is executable rather than
        merely optimal.
  \item The data handling is justified operation by operation (fill, truncate,
        interpolate), and the interpolation error is bounded by the sampling step.
  \item The sensitivity analysis is run per problem, and the stiffness peak is explained
        by the natural-frequency crossing rather than reported as a number.
\end{enumerate}

\subsection{Weaknesses}

% 不足要写"会因为什么原因、影响哪个结论"，最好带一条改进方向。
\begin{enumerate}
  \item Parameters are treated as constant; load, wear and temperature make $M$, $k$,
        $c$ time-varying, so the robustness claim holds only near the nominal set.
  \item Friction in the eccentric-block bearings and losses in the bevel-gear train are
        ignored, which biases the derived angular acceleration and hence $F_a$.
  \item The observer is tuned on the nominal model; gain scheduling for a strongly
        time-varying disturbance is left open.
  \item Only a single disturbance realisation per scenario is reported; a Monte Carlo
        sweep over disturbance phase would give the suppression rate a confidence band.
\end{enumerate}

% ┈┈┈┈┈┈┈┈┈┈┈┈┈┈┈┈┈┈┈┈┈┈┈┈┈┈┈┈┈┈┈┈┈┈┈
% 【Section writing prompt】
% ════════════════════════════════════════════
% Counts:   at least 4 strengths and 4 weaknesses. Three-and-two (the usual) reads as
%           filler.
% Every item names a mechanism and its consequence: "friction in the bearings biases the
%           angular acceleration, hence F_a" beats "the model is simplified".
% Weakness items should hint at the fix — that is what turns a weakness list into a
%           research direction.
% ════════════════════════════════════════════
